% -*- coding: utf-8; mode: latex; -*-

% +++
% latex = "uplatex"
% bibtex = "upbibtex"
% +++

%
% FIT2024 向け LaTeX クラスファイル
% https://github.com/trueroad/FITpaper-class
%

\documentclass[dvipdfmx]{FITpaper}

\usepackage{graphicx}
\usepackage{dblfloatfix}
\usepackage{multirow}
\usepackage{booktabs}
\usepackage{flushend}

\makeatletter
\renewcommand{\ackbibsection}[1]{%
  \par
  \vspace{6bp}%
  {\centering\normalsize\sffamily #1\par}%
  \vspace{2pt}%
}
\makeatother

% 和文タイトル
\jtitle{空き家マッチングにおける借り手の熱意を代弁する\\交渉対話エージェントの試作}

% 欧文タイトル
\etitle{Prototyping a Negotiation Dialogue Agent as a Proxy for Tenant Enthusiasm in Vacant House Matching}

% 著者数：著者の数だけ c を書く
\authors{cccc}

% 和文著者名
\jauthors{%
  福富 隆大\affmark{1} &
  山崎 紘介\affmark{1} &
  長澤 史記\affmark{1} &
  白松 俊  \affmark{1}
}

% 欧文著者名
\eauthors{%
  Takahiro Fukutomi &
  Kosuke Yamazaki   &
  Fumiki Nagasawa   &
  Shun Shiramatsu
}

% 所属
\afftext{1}{名古屋工業大学  Nagoya Institute of Technology.}


\begin{document}

\maketitle

% ======================================
\section{はじめに}
% ======================================

日本では空き家の増加が社会問題となっており，その活用が求められている．
総務省の住宅・土地統計調査（2023年）によると，空き家数は約900万戸に達し過去最多を更新した．
しかし，空き家オーナー（大家）の多くは積極的に物件を貸し出そうとしているわけではなく，
使い方や借り手の人柄によっては貸してもよいと考えている層が存在する．

このような消極的なオーナー層にはたらきかけるマッチングサービスに「さかさま不動産」がある．
従来の不動産取引では大家が借り手を選ぶのに対し，さかさま不動産では借り手が自身の活動目的や
想いを記事として公開し，大家が「この人に貸したい」と思う相手を選ぶ．
借り手は自身の夢・目的・背景を記事として公開しており，
大家がその記事を読んで関心を持った場合に対話が始まる．
ただし，大家は対話時点で記事を詳しく読んでいない場合も多く，
借り手が自ら自己紹介・目的説明・熱意のアピールを行う必要がある．

こうした「申請者が選定者に熱意を伝えながら適合性を確認する対話」は，
さかさま不動産に限らず，就職活動における志望者の自己アピール，
研究資金申請における研究者と評価委員のやりとりなど，
様々な場面に共通する課題構造である．
すなわち，\textbf{情報の非対称性（相手が申請者の詳細を事前に知らない）}のもとで
申請者が熱意と適合性の両方を伝えることが求められるという課題は，広く一般に適用可能な問題として捉えられる．
多数の相手に対して自ら行うことは負担が大きく，不慣れな利用者には難しい．

本研究は，こうした状況において申請者（借り手）の記事をもとに交渉対話を代行する
AI エージェントを試作することを最終的な目標とする．
これにより，不慣れな利用者でも大家に熱意を届けられる機会を提供するとともに，
感情表現が対話品質に与える影響を実証的に明らかにする価値がある．
そのアプローチとして，感情表現の指示をプロンプトに追加する条件と追加しない条件を比較する
予備実験を行い，システムの実現可能性と評価手法を検討する．

% ======================================
\section{関連研究}
% ======================================

\subsection{ロールプレイング言語エージェント（RPLA）}

LLM を用いてペルソナを与え，特定の人物として対話させるロールプレイング言語エージェント
（RPLA）の研究が進んでいる~\cite{chen2024rpla,tseng2024twotales}．
RPLA における課題として，LLM がキャラクターとして与えられた情報の範囲を超えた
事実を発話する hallucination が挙げられている~\cite{yu2024timechara}．
本研究が対象とする借り手記事は，特定個人の目的・価値観を記述した
Individualized Persona~\cite{chen2024rpla}に近く，
記事にない事実を発話するリスクへの対処として，捏造防止ルールをプロンプトに組み込んだ．

\subsection{LLMエージェントにおける感情表現}

LLM エージェントの感情表現に関して，
進化的強化学習により感情ポリシーを動的に最適化する EvoEmo~\cite{evoemo2025}，
感情状態を記憶検索に組み込む Emotional RAG~\cite{emotionalrag2024}，
戦略的計画と表現最適化を組み合わせた DMNA~\cite{dmna2025} などが提案されている．
これらの研究は，感情表現が LLM エージェントの説得力やペルソナ整合性に寄与することを示している．
本研究ではこれらの知見を踏まえ，感情表現指示プロンプティングによる簡易的な実現を試みる実験を行った．


% ======================================
\section{提案システム}
% ======================================

\subsection{システム概要}

提案システムは，借り手の記事を入力として大家との多ターン対話を行う
借り手エージェントである．
借り手エージェントは主に2つの機能を担う．

\begin{enumerate}
  \item \textbf{感情表現（主）}：借り手の動機を率直に伝え「この人に貸したい」と感じさせる発話を生成する．
  \item \textbf{適合確認の質問（従）}：物件が借り手の目的に合うか確認する質問を行う（2〜3回程度）．
\end{enumerate}

感情表現の効果を定量的に評価することが本研究の主目的である．

\subsection{借り手エージェントの設計}

エージェントには借り手記事の全文と発話指示をシステムプロンプトとして与える．
2条件の違いは発話指示の内容のみであり（表\ref{tab:condition}），
それ以外の借り手情報・捏造防止ルールは両条件で同一とした．

\begin{table}[t]
  \centering
  \caption{条件の比較}
  \label{tab:condition}

  \begin{tabular}{lp{4.5cm}}
    \toprule
    条件 & 発話指示の内容 \\
    \midrule
    Baseline & 借り手の目的・希望を事実に基づき伝える．感情表現指示なし． \\
    Proposed & Baseline に加え，フェーズ別の感情表現指示を付与．opening では動機を率直に，middle では懸念を受け止めてから質問，closing では前向きな姿勢を示す． \\
    \bottomrule
  \end{tabular}
\end{table}

捏造防止のため，記事に根拠のない断定的表現を避けるルールを
両条件共通でプロンプトに組み込んだ．

% ======================================
\section{実験}
% ======================================

\subsection{実験設定}

提案システムの2条件（Baseline・Proposed）を ChatGPT の Custom GPTs として実装し，
人間の参加者が大家役を担う対話実験を行った．
感情表現の指示の有無のみを独立変数とし，同じ借り手情報から生成される発話の品質差を比較する．

互いに異なる借り手ケースを割り当てた6種類のGPTsを用意した（表\ref{tab:gpts}）．
参加者（研究室メンバー6名）は物件情報シートを参照しながら大家役として各GPTと4ターン程度対話し，
全6 GPTsを評価した．
なお，1つの借り手記事に対して Baseline または Proposed のいずれか一方の GPT を割り当てた．

\begin{table}[t]
  \centering
  \caption{実験に使用した借り手エージェント GPTs}
  \label{tab:gpts}
  \small
  \begin{tabular}{llp{1.55cm}}
    \toprule
    GPT名 & 借り手 & 条件 \\
    \midrule
    整体院 & パーソナルトレーナー・整体師 & Baseline \\
    託児カフェ & デザイナー・ベーグル屋 & Baseline \\
    陶作家 & 陶作家 & Baseline \\
    アウトドア & キャンプインストラクター & Proposed \\
    ピラティス & ピラティス講師・管理栄養士 & Proposed \\
    学童クラブ & 学童クラブ運営者 & Proposed \\
    \bottomrule
  \end{tabular}
\end{table}

\subsection{評価軸と評価方法}

表\ref{tab:survey}にアンケート項目を示す．リカートスケールは7点を最良・1点を最低とする．

\begin{table}[h]
  \centering
  \caption{アンケート項目}
  \label{tab:survey}
  \begin{tabular}{lp{3.2cm}l}
    \toprule
    項目 & 質問文（要約） & 尺度 \\
    \midrule
    根拠忠実性 & 記事にない事実を言っていた件数 & 違反件数 \\
    熱意の伝達 & 借り手の想いが伝わった & 7件法 \\
    人となり & どんな人かわかった & 7件法 \\
    使い方 & この物件をどう使うか想像できた & 7件法 \\
    貸したいか & 次の話に進みたいと思う & 7件法 \\
    \bottomrule
  \end{tabular}
\end{table}

% ======================================
\section{結果と考察}
% ======================================

\subsection{根拠忠実性}

全36件（6 GPTs $\times$ 6名）の対話において違反件数は0件であった．
ただし，計算から導いた数値（例：「35年以上続いてきた」を1991年創設から算出）を
用いた発話は観察された．これらは捏造ではなく，算術的に正当な推論として許容した．

\subsection{条件間の比較}

各 GPT の評価平均を表\ref{tab:results_gpt}に，条件別平均を表\ref{tab:results_cond}に示す．

\begin{table}[h]
  \centering
  \caption{GPT別評価平均（n=6）}
  \label{tab:results_gpt}
  \small
  \begin{tabular}{llcccc}
    \toprule
    GPT & 条件 & 熱意 & 人となり & 使い方 & 貸したい \\
    \midrule
    整体院 & B & 4.33 & 3.83 & 5.83 & 4.83 \\
    託児カフェ & B & 5.50 & 5.17 & 6.17 & 5.50 \\
    陶作家 & B & 5.17 & 4.67 & 6.33 & 4.67 \\
    アウトドア & P & 5.00 & 3.83 & 5.00 & 4.50 \\
    ピラティス & P & 5.67 & 5.00 & 5.50 & 5.33 \\
    学童クラブ & P & 5.83 & 4.00 & 5.00 & 5.17 \\
    \bottomrule
  \end{tabular}
\end{table}

\begin{table}[h]
  \centering
  \caption{条件別評価平均（Mann-Whitney U, $n=18$/条件）}
  \label{tab:results_cond}
  \begin{tabular}{p{2.0cm}cccc}
    \toprule
    条件 & 熱意 & 人となり & 使い方 & 貸したい \\
    \midrule
    Baseline (B) & 5.00 & 4.56 & \textbf{6.11} & 5.00 \\
    Proposed (P) & \textbf{5.50} & 4.28 & 5.17 & 5.00 \\
    \midrule
    $p$値 & .143 & .807 & \textbf{.002} & .629 \\
    \bottomrule
  \end{tabular}
\end{table}

\subsection{考察}

Mann-Whitney U 検定（両側，$n=18$/条件）の結果，
使い方のイメージスコアで Baseline が Proposed を有意に上回った
（$U=253.5, p=.002$）．
感情表現指示の付加により感情的訴求に重心が移り，
物件利用イメージの提示が相対的に少なくなった結果と考えられる．

熱意の伝達スコアは Proposed が Baseline をやや上回ったが（5.50 vs 5.00），
有意差は認められなかった（$U=118.5, p=.143$）．
感情表現指示の効果を確認するには評価者数・GPT数の拡充が必要である．

人となり（$p=.807$）および貸したいと思うか（$p=.629$）は
両条件で有意差を示さなかった．

なお，各評価者が複数 GPT を評価しているため評価値は完全独立ではなく，
検定の解釈には留意が必要である．
また，各条件に異なる借り手ケースを割り当てたため，
条件間の差異は感情表現指示の効果と借り手ケースの特性が交絡した参考値にとどまる
（需要特性バイアス回避のためこの設計を採用した）．

評価者からの定性的コメントとして，
「AIらしい文章」「話が長い」「身の上話後の質問で話題がずれる」
「答えにくい地域質問（地域との交流の雰囲気など）への対応が苦手」
といった観察が得られた．これらは GPTs の制約（多段 LLM 呼び出しが不可なため計画生成を単一プロンプトで代替したこと，ターン数に制限があること）に起因すると考えられる．


% ======================================
\section{おわりに}
% ======================================

本研究では，借り手の熱意を代弁して大家と交渉対話を行う AI エージェントを試作し，
感情表現指示の有無を独立変数とした比較実験を行った．
根拠忠実性の違反は全36件で0件であり，捏造防止の観点では有効に機能した．
使い方のイメージスコアでは Baseline が Proposed を有意に上回り（$p=.002$），
熱意の伝達では Proposed がやや上回ったが有意差には至らなかった（$p=.143$）．
感情表現指示が発話の重心を変える可能性が示唆された．

今後の課題として，交絡の解消（同一ケースを被験者間でカウンターバランスし，
各参加者は一方の条件のみ評価する計画）および評価者数の拡大が挙げられる．

% ======================================
\acknowledgment{%
  本研究の一部はJSPS科研費（24K03052, 25K00866）の支援を受けた．
}
% ======================================

\bibliographystyle{junsrt}
\bibliography{references}

\end{document}
